\begin{abstract}
We establish the quasi-Riemann hypothesis by proving that every Dirichlet $L$-function, including Riemann's zeta function, has no zeros in the half-plane $\Rea s>11/12$. More generally, we prove the same zero-free half-plane for every finite-order Hecke $L$-function over $K=\Q(\sqrt{-3})$. In particular, this rules out the existence of Landau--Siegel zeros.
\end{abstract}

\tableofcontents

\section{Introduction}\label{sec:intro}

For a primitive Dirichlet character $\chi$ of conductor $q$, the associated \emph{Dirichlet $L$-function} is defined for $\Rea s>1$ by the formula
\[
 L(s,\chi)=\sum_{n\ge1}\frac{\chi(n)}{n^s}, \qquad \Rea s>1,
\]
and admits a meromorphic continuation to $s \in \C$ \cite[\S4.6]{IK04}. When $\chi$ is the trivial character, $L(s, \chi)$ recovers the Riemann zeta function $\zeta(s)$ \cite{Rie59}.

The zeros of Dirichlet $L$-functions are of significant interest, such as for their role in governing the distribution of primes in arithmetic progressions. The \emph{Generalized Riemann Hypothesis} predicts that all zeros of $L(s, \chi)$ in the critical strip $0 < \Rea s < 1$ satisfy $\Rea s = 1/2$.
For $\zeta(s)$, the weaker assertion that there exists
$\varepsilon>0$ such that $\zeta(s)$ has no zeros in
$\Rea s>1-\varepsilon$ is called a \emph{quasi-Riemann Hypothesis}
(e.g., in \cite[\S1]{BG17}, \cite[p.~274]{MS02}, and \cite[\S2]{BR18}).
For Dirichlet $L$-functions, we consider the analogous assertion
with a single $\varepsilon>0$ valid for every primitive character
$\chi$, independently of both conductor and height.



Let $K=\Q(\sqrt{-3})$. The main result of this paper is the following.

\begin{theorem}\label{thm:main}
Every finite-order Hecke $L$-function over $K$ has no zeros in the half-plane $\Rea s>11/12$. Thus every Dirichlet $L$-function, including $\zeta(s)$, has no zeros for $\Rea s>11/12$.
\end{theorem}
